\documentclass[12pt]{article}

\usepackage[margin=1in]{geometry}
\usepackage{amsmath}
\usepackage{array}
\usepackage{enumitem}
\usepackage{hyperref}

\setlength{\parindent}{0pt}
\setlength{\parskip}{0.8em}

\begin{document}

\section*{Busy Time Scheduling Project: Final Project Summary}

\section*{Project Repository and File Structure}

The work for the project is being organized in the GitHub repository

\[
\texttt{https://github.com/pacots/busy-time}
\]

This repository contains the main Python scripts, input instances, output files,
papers, and notes used throughout the project. The current structure is organized
roughly as follows:

\begin{itemize}
    \item \texttt{inputs/}: contains the CSV input files used for testing the
    algorithms. This includes the original small examples such as
    \texttt{input1.csv} through \texttt{input5.csv}, as well as larger generated
    instances such as 50-job, 99-job, 100-job, 150-job, 200-job, and 1000-job
    cases.

    \item \texttt{outputs/}: contains the output files generated after running
    the scheduling scripts. These outputs are used to compare schedules, busy
    time values, and feasibility results.

    \item \texttt{papers/}: contains the paper "LP rounding and combinatorial algorithms for minimizing active and busy time", which contains Theorem 7.

    \item \texttt{docs/}: contains documentation-related files for the project.

    \item \texttt{meeting-notes/}: contains notes from the group meetings and
    project discussions.

    \item \texttt{generate\_job\_instances.py}: script used to generate new job
    instances and save them as CSV files.

    \item \texttt{theorem7.py} and
    \texttt{theorem7\_scheduler\_separate\_outputs.py}: previous scripts related
    to the Theorem 7 preemptive busy time scheduling approach. These were useful
    earlier in the project, but they are not the main scripts currently being
    used.

    \item \texttt{l2\_preemptive\_busy\_time.py}: script for the greedy/preemptive
    busy time approach that we have been testing from ”LP rounding and combinatorial algorithms for minimizing active and busy time” Theorem 7. It is an approximating algorithm and produces feasible solutions, but not necessarily optimum.

    \item \texttt{local\_search.py} and \texttt{local\_search\_busy\_time.py}:
    scripts related to the local search / greedy improvement approach. They are failed attempts for the implementation of the local search algorithm from spaa18.pdf "Brief Announcement: A greedy 2 approximation for the active time problem", Appendix A.3.

    \item \texttt{ilp\_busy\_time\_pulp\_gurobi.py}: script for solving the busy
    time problem using an ILP formulation with PuLP and Gurobi. This program is AI generated by ChatGPT. We observed that it is able to solve for 100 job instances, but the runtime becomes too large for 150 job instances, and did not finish within a 10 minute window.

    \item \texttt{ilp\_busy\_time\_pulp\_gurobi\_time\_limit.py}: version of the
    ILP script that has not yet been tested. It prints the best solution obtained within a certain time.

    \item \texttt{check\_feasibility.py}: script used to check whether a produced
    schedule is feasible.
\end{itemize}

The main files in use are:

\begin{itemize}[leftmargin=*]
    \item \texttt{generate\_job\_instances.py}
    \item \texttt{l2\_preemptive\_busy\_time.py}
    \item \texttt{ilp\_busy\_time\_pulp\_gurobi.py}
    \item \texttt{ilp\_busy\_time\_pulp\_gurobi\_time\_limit.py}
    \item \texttt{check\_feasibility.py}
\end{itemize}

Together, these scripts let us generate test inputs, run a greedy/preemptive
scheduling approach, run an ILP-based approach, and check whether the schedules
that are produced are feasible. The use of the ILP scripts requires Gurobi. A
free academic license can be obtained by registering with the school's email and
generating the license while connected to campus WiFi.

\section*{Input File Format}

The scheduling scripts expect CSV files in the format used by the project
inputs. The first rows give the problem type and the machine capacity, followed
by the job table. For example:

\begin{verbatim}
Busy Time
Capacity
2
Job,Release,Deadline,processingTime
A,1,4,2
B,1,6,3
C,2,5,1
\end{verbatim}

Each job row contains the job name, release time, deadline, and processing time.
The processing time must fit inside the job's window, so each job should satisfy

\[
\text{processing time} \leq \text{deadline} - \text{release time}.
\]

Most input files are stored in the \texttt{inputs/} folder.

\section*{Output File Organization}

The main scheduling scripts save their results in the \texttt{outputs/} folder.
Each run usually creates a separate output folder named using the input file and
the algorithm. Inside that folder, the scripts save CSV files for the schedule
and summary results.

For the greedy/preemptive script l2 preemptive busy time.py, the output folder usually contains files such
as:

\begin{itemize}[leftmargin=*]
    \item \texttt{\_input\_jobs.csv}: cleaned copy of the input jobs.
    \item \texttt{\_active\_iterations.csv}: record of the active intervals added
    during the greedy step.
    \item \texttt{\_unbounded\_active\_intervals.csv}: active intervals from the
    unbounded-capacity schedule.
    \item \texttt{\_unbounded\_schedule.csv}: preemptive job schedule before
    applying the machine capacity limit.
    \item \texttt{\_bounded\_schedule.csv}: final bounded schedule after applying
    capacity \(g\).
    \item \texttt{\_summary.csv}: summary values such as capacity, unbounded busy
    time, bounded busy time, and ratio.
\end{itemize}

For the ILP script ilp busy time pulp gurobi.py, the output folder usually contains files such as:

\begin{itemize}[leftmargin=*]
    \item \texttt{\_input\_jobs.csv}: cleaned copy of the input jobs.
    \item \texttt{\_bounded\_schedule.csv}: final schedule produced by the ILP.
    \item \texttt{\_job\_schedule.csv}: job-by-job version of the ILP schedule.
    \item \texttt{\_summary.csv}: solver status, capacity, candidate machines,
    and bounded busy time.
\end{itemize}

The feasibility checker can be run on these output folders or directly on a
specific bounded schedule CSV file.

\section*{1. \texttt{generate\_job\_instances.py}}

The purpose of \texttt{generate\_job\_instances.py} is to create new CSV input
files for the busy time scheduling problem. This is useful because the original
examples are small, and we need more test cases in order to see how the
algorithms behave on different kinds of inputs.

The generated CSV files follow the same general format as the project inputs.
Each job has:

\begin{itemize}[leftmargin=*]
    \item a job name,
    \item a release time,
    \item a deadline,
    \item a processing time.
\end{itemize}

The script makes sure that the generated jobs are valid. In particular, each job
must have enough room to fit inside its own time window. In other words, the
processing time cannot be larger than the difference between the deadline and the
release time:

\[
\text{processing time} \leq \text{deadline} - \text{release time}.
\]

The script can generate different types of instances. For example, some cases are
meant to have lower overlap between jobs, while others are meant to have more
overlap, tighter windows, looser windows, heavier processing times, lighter
processing times, clustered overlap, nested windows, or more random behavior.
This gives us a way to test the scheduling scripts on more than one style of
problem.

The script can be run from the command line. For example, one possible command is:

\begin{verbatim}
python generate_job_instances.py --num-jobs 50 --case-type low_overlap \
    --capacity 2 --output-dir inputs
\end{verbatim}

There is also an option to generate several case types at once and save them in
an output folder:

\begin{verbatim}
python generate_job_instances.py --num-jobs 50 --all-cases \
    --capacity 2 --output-dir inputs
\end{verbatim}

\section*{2. \texttt{l2\_preemptive\_busy\_time.py}}

The file \texttt{l2\_preemptive\_busy\_time.py} is the script for the
greedy/preemptive scheduling approach. This is the faster scheduling script we
have been using to produce schedules without solving a full optimization model.

The script reads an input CSV file, stores the jobs, and then builds a schedule
in two main steps.

First, it creates an unbounded-capacity schedule. In this step, the capacity
limit is temporarily ignored. The script looks at the jobs and opens active time
intervals where jobs can be placed. The general idea is to make sure each job has
enough active time inside its release-deadline window. The script tends to place
work as late as possible within the allowed windows.

Second, it converts that schedule into a bounded-capacity schedule. This is where
the machine capacity \(g\) is used. If more than \(g\) jobs are running at the
same time, the script splits them across multiple machines. For example, if
\(g=2\), then each machine can handle at most two jobs at once.

The script produces several output files, including:

\begin{itemize}[leftmargin=*]
    \item a cleaned version of the input jobs,
    \item the active intervals from the unbounded part,
    \item the unbounded schedule,
    \item the final bounded schedule,
    \item a summary file with values such as capacity, unbounded busy time,
    bounded busy time, and ratio.
\end{itemize}

A typical command is:

\begin{verbatim}
python l2_preemptive_busy_time.py input1.csv
\end{verbatim}

The script looks for the input file in the \texttt{inputs/} folder first. The
results are written to the \texttt{outputs/} folder.

\section*{3. \texttt{ilp\_busy\_time\_pulp\_gurobi.py}}

The file \texttt{ilp\_busy\_time\_pulp\_gurobi.py} is the script for the ILP
approach. The main purpose of this script is to model the busy time scheduling
problem as an optimization problem and solve it using a solver such as Gurobi
through PuLP.

Instead of using a greedy rule, the ILP script gives the solver the scheduling
rules as constraints. The solver then tries to find a schedule with the smallest
possible busy time.

At a high level, the ILP model keeps track of:

\begin{itemize}[leftmargin=*]
    \item whether a job is assigned to a machine at a certain time,
    \item whether a machine is active at a certain time,
    \item whether each job gets all of its required processing time,
    \item whether the capacity limit is respected.
\end{itemize}

The objective of the ILP is to minimize total busy time. This means the solver is
trying to use as little total machine activity as possible while still scheduling
every job correctly.

The benefit of the ILP approach is that it is more optimization-based than the
greedy script, so it is useful for comparison. However, the downside is that it
can become slow. For some instances, the Gurobi search may take an unreasonable
amount of time, especially when the input is larger or more complicated.

A typical command is:

\begin{verbatim}
python ilp_busy_time_pulp_gurobi.py input1.csv
\end{verbatim}

A time limit can also be passed to the solver:

\begin{verbatim}
python ilp_busy_time_pulp_gurobi.py input1.csv --time-limit 300
\end{verbatim}

This tells the solver to stop after 300 seconds. The time-limit version of the
script is also being worked on because we want to make sure that useful output is
saved even when the solver does not fully finish.

\section*{4. \texttt{check\_feasibility.py}}

The file \texttt{check\_feasibility.py} is used to check whether a schedule is
actually valid. This is important because after a schedule is produced, we still
need to confirm that it follows the rules of the problem.

The checker reads the original input file and a bounded schedule file. It then
checks basic feasibility conditions, such as:

\begin{itemize}[leftmargin=*]
    \item every scheduled job exists in the input,
    \item every job is scheduled only inside its release-deadline window,
    \item every job receives the correct total amount of processing time,
    \item a job is not running in two places at the same time,
    \item no machine is assigned more than \(g\) jobs at once,
    \item the schedule file has the expected format.
\end{itemize}

This script is useful because it separates the act of producing a schedule from
the act of checking a schedule. Even if the greedy script or ILP script runs
successfully, the feasibility checker gives us another way to confirm that the
output makes sense.

The checker can be run using an output directory:

\begin{verbatim}
python check_feasibility.py inputs/input1.csv outputs/input1_l2_preemptive
\end{verbatim}

It can also be run directly on a bounded schedule file:

\begin{verbatim}
python check_feasibility.py inputs/input1.csv \
    outputs/input1_l2_preemptive/input1_l2_preemptive_bounded_schedule.csv
\end{verbatim}

\section*{Current Status and Next Step}

At this stage, we have a basic workflow for the project:

\begin{enumerate}[leftmargin=*]
    \item Generate job instances with \texttt{generate\_job\_instances.py}.
    \item Run the greedy/preemptive scheduler with
    \texttt{l2\_preemptive\_busy\_time.py}.
    \item Run the ILP scheduler with
    \texttt{ilp\_busy\_time\_pulp\_gurobi.py}.
    \item Check the resulting schedules with \texttt{check\_feasibility.py}.
    
\end{enumerate}

The main problem we are currently seeing is that the ILP/Gurobi approach can
sometimes take too long. Because of that, one future goal is to make the solver
part more practical by adding a time limit and saving the best solution found so
far when the time limit is reached.

\begin{enumerate}[leftmargin=*]
    \item Text and fix if needed ILP version with time limit ilp busy time pulp gurobi time limit.py.
    \item Make local search from Appendix A.3 work
    \item Investigate if a dynamic program exists for g=2.
\end{enumerate}

\end{document}